\section{Historial de cambios y compatibilidad}
Aquí reúno los cambios que han dado forma al núcleo actual y las decisiones
que conservo para seguir leyendo textos anteriores.

\subsection{Ampliación del núcleo}
\begin{itemize}
\item Incorporé veinte enunciados T01--T20 con contexto, lectura esperada y
contraste. Conservo los diez ejemplos iniciales E01--E10.
\item Añadí la selección por condición para componer conceptos, junto con
cinco conceptos base y tres propiedades contextuales.
\item Amplié el catálogo a nueve relaciones, con pertenencia, influencia causal
en un modelo, gusto, satisfacción de propiedad y norte geográfico.
\item Incorporé alternativa, posibilidad, pregunta, selección, secuencia y
mientras; el inventario reúne ahora veintidós signos.
\item Definí colecciones enumeradas, asignación y reglas de tipos, alcance y contexto.
\item Adopté las formas simplificadas de Colección y Fuente y ajusté las
uniones de Mientras y Espacio.
\end{itemize}

\subsection{Lectura de textos anteriores}
Conservo las funciones de K01--K16. Los cambios de dibujo de Colección,
Fuente y Espacio mantienen sus significados. Los usos indexados de K01,
K02, K04 y K06 añaden construcciones que antes no estaban definidas.

Mantengo disponibles las macros anteriores. Archivo las gramáticas previas en
\path{anexos/gramatica-inicial.tex} y \path{anexos/extension-anterior.tex};
reúno las reglas vigentes en \path{docs/gramatica.tex}.

Uso las formas B como grafías habituales de K06/K16. Conservo V06/V16 como
alias de edición y H06/H16 para representar las formas anteriores. Estos
códigos auxiliares no añaden significados al catálogo.

\subsection{Cómo interpreto las comprobaciones}
Distingo la compilación, que comprueba que los documentos pueden representarse,
de las pruebas semánticas, que verifican contrastes concretos. Las lecturas
esperadas acompañan a las expresiones como claves de estudio; no las considero
una prueba independiente de comprensión.

Uso el cuaderno para anotar mis dificultades con signos, referencias y alcance.
Quiero que las siguientes revisiones partan de esas observaciones.

\subsection{Consolidación del manual}
Reuní la gramática y el trazado de los 22 signos en el cuerpo de referencia y
separé el historial en anexos. Añadí las hojas \path{lectura.tex} y
\path{relectura.tex}, junto con M01, para practicar expresiones completas.
También precisé la continuidad de referencias entre páginas y el uso de
nombres definidos dentro de un texto.

Conservo respaldos separados de las etapas anteriores. Los sufijos numéricos
de algunos archivos son referencias internas; el nombre de mi escritura
sigue siendo KalaScript.
